\documentclass[aspectratio=169, 14pt]{beamer}
\usepackage{stampfly_slides}

\lesson{7}{テレメトリ + ステップ応答}{Telemetry + Step Response}

\begin{document}

% ------------------------------------------------------------------
\begin{frame}
  \titlepage
\end{frame}

% ------------------------------------------------------------------
\begin{frame}{今日のゴール / Today's Goal}
  \begin{block}{目標}
    WiFi テレメトリでステップ応答データを取得し、制御性能を分析する
  \end{block}

  \vspace{1em}

  \begin{itemize}
    \item ステップ応答とは何か
    \item 立ち上がり時間・オーバーシュート・整定時間
    \item \code{sf log wifi} でデータ取得
  \end{itemize}
\end{frame}

% ------------------------------------------------------------------
\begin{frame}{ステップ応答 / Step Response}
  \begin{center}
    \resizebox{!}{0.78\textheight}{\documentclass[border=10pt]{standalone}
\usepackage{tikz}
\usetikzlibrary{arrows.meta}

\begin{document}

\definecolor{refcolor}{RGB}{0, 120, 200}    % Blue - reference
\definecolor{respcolor}{RGB}{220, 50, 30}   % Red - response
\definecolor{gridcolor}{RGB}{200, 200, 200}
\definecolor{labelcolor}{RGB}{80, 80, 80}
\definecolor{annocolor}{RGB}{40, 160, 40}   % Green - annotations
\begin{tikzpicture}[
    axis/.style={-{Stealth[length=3mm]}, line width=1pt, labelcolor},
    ref/.style={refcolor, line width=2pt},
    resp/.style={respcolor, line width=2pt},
    anno/.style={annocolor, line width=1pt, dashed},
    dim/.style={{Stealth[length=2mm]}-{Stealth[length=2mm]}, annocolor, line width=0.8pt},
    annolbl/.style={font=\small\bfseries, annocolor},
  ]

  % Grid
  \foreach \x in {0, 1, ..., 12} {
    \draw[gridcolor] (\x, 0) -- (\x, 6);
  }
  \foreach \y in {0, 1, ..., 6} {
    \draw[gridcolor] (0, \y) -- (12, \y);
  }

  % Axes
  \draw[axis] (-0.3, 0) -- (12.5, 0) node[right, font=\bfseries] {$t$ [s]};
  \draw[axis] (0, -0.3) -- (0, 6.5) node[above, font=\bfseries] {$\omega$ [rad/s]};

  % Time labels
  \foreach \x/\t in {0/0, 2/0.5, 4/1.0, 6/1.5, 8/2.0, 10/2.5, 12/3.0} {
    \node[font=\scriptsize, below, labelcolor] at (\x, -0.1) {\t};
  }

  % Reference (step input at t=0.5s = x=2)
  \draw[ref] (0, 0) -- (2, 0) -- (2, 4) -- (12, 4);
  \node[font=\small\bfseries, refcolor, above left] at (1, 0.3) {Target};

  % Response curve — exact second-order underdamped step response
  % y(τ) = A · [1 − exp(−ζωₙτ) / √(1−ζ²) · sin(ωd·τ + φ)]
  % ζ = 0.3, ωₙ = 6.5 rad/s, ωd = 6.20 rad/s, φ = arccos(ζ), A = 4.0
  \pgfmathsetmacro{\zetawn}{1.95}       % ζωₙ = 0.3 × 6.5
  \pgfmathsetmacro{\wddeg}{355.24}      % ωd [deg/s] = 6.20 × 180/π
  \pgfmathsetmacro{\phideg}{acos(0.3)}  % φ = arccos(ζ) [deg]
  \pgfmathsetmacro{\invsqrt}{1/sqrt(1-0.09)} % 1/√(1−ζ²)
  \draw[resp] (0, 0) -- (2, 0);
  \draw[resp] plot[domain=0:2.5, samples=200]
    ({2 + \x*4}, {4*(1 - \invsqrt*exp(-\zetawn*\x)*sin(\wddeg*\x + \phideg))});
  \node[font=\small\bfseries, respcolor, above right] at (9, 4.1) {Response};

  % --- Annotations ---

  % Overshoot (peak: τ=π/ωd=0.507s → x=4.03, y=5.49)
  \draw[anno] (4.0, 4) -- (4.0, 5.5);
  \draw[dim] (4.2, 4) -- (4.2, 5.5);
  \node[annolbl, right] at (4.3, 4.75) {Overshoot};

  % Rise time
  \draw[anno] (2, -0.8) -- (2, -0.3);
  \draw[anno] (3.2, -0.8) -- (3.2, -0.3);
  \draw[dim] (2, -0.6) -- (3.2, -0.6);
  \node[annolbl, below] at (2.6, -0.85) {Rise Time};

  % Settling time
  \draw[anno] (2, -1.8) -- (2, -1.5);
  \draw[anno] (8.2, -1.8) -- (8.2, -1.5);
  \draw[dim] (2, -1.65) -- (8.2, -1.65);
  \node[annolbl, below] at (5, -1.85) {Settling Time};

  % Steady-state value line
  \draw[refcolor, line width=0.5pt, dashed] (0, 4) -- (12, 4);

  % Settling band (5%)
  \fill[annocolor, opacity=0.08] (0, 3.8) rectangle (12, 4.2);
  \node[font=\scriptsize, annocolor, right] at (10, 3.2) {$\pm 5\%$ band};

\end{tikzpicture}
\end{document}
}
  \end{center}
\end{frame}

% ------------------------------------------------------------------
\begin{frame}{テレメトリと実験 API}
  \begin{block}{WiFi テレメトリ（自動送信）}
    システムが IMU・姿勢・センサデータを 400\,Hz で自動送信 \\
    \code{sf log wifi} で受信 → CSV 保存 → \code{sf log viz} で可視化
  \end{block}

  \vspace{0.3em}

  \begin{table}
    \centering
    \begin{tabular}{lll}
      \hline
      \textbf{関数} & \textbf{説明} & \textbf{引数} \\
      \hline
      \api{led\_color(r, g, b)} & LED 色設定（フェーズ表示用） & 各 0--255 \\
      \hline
    \end{tabular}
  \end{table}

  \vspace{0.3em}

  \begin{block}{実験手順}
    \begin{enumerate}
      \item ホバリング状態でステップ入力を自動付加
      \item \code{sf log wifi} でリアルタイム記録
      \item CSV を分析して性能指標を計測
    \end{enumerate}
  \end{block}
\end{frame}

% ------------------------------------------------------------------
\begin{frame}[fragile]{実習: ステップ応答実験 / Hands-on}
  \begin{lstlisting}[style=sfcpp, title={student.cpp (excerpt)}]
// Step response: 2s wait -> 1s step -> 1s recovery
float step_rate = 0.5f;   // Step amplitude [rad/s]
uint32_t delay  = 800;    // 2s at 400Hz
uint32_t dur    = 400;    // 1s step
float roll_step = 0.0f;
if (elapsed >= delay && elapsed < delay + dur)
    roll_step = step_rate;  // Step ON!
float roll_target = roll_step;  // Override stick
// LED: red during step, green otherwise
if (roll_step > 0) ws::led_color(50, 0, 0);
else               ws::led_color(0, 50, 0);
// gyro_x = roll rate is auto-sent via WiFi telemetry
  \end{lstlisting}
\end{frame}

% ------------------------------------------------------------------
\begin{frame}{チェックポイント / Checkpoint}
  \begin{alertblock}{確認事項}
    \begin{itemize}
      \item[$\square$] \code{sf log wifi} でデータを取得できた
      \item[$\square$] 波形でオーバーシュートを確認
      \item[$\square$] 整定時間を計測（目安: 0.3s 以内）
    \end{itemize}
  \end{alertblock}

  \vspace{1em}

  \begin{block}{次のレッスン}
    Lesson 8: モデリング --- 運動方程式からミキサーを導出
  \end{block}
\end{frame}

\end{document}
